\documentclass[11pt]{article}

\usepackage[margin=1in]{geometry}
\usepackage{amsmath}
\usepackage{booktabs}
\usepackage{float}
\usepackage{graphicx}
\usepackage{hyperref}
\usepackage{natbib}
\usepackage{url}
\Urlmuskip=0mu plus 1mu\relax

\title{Fixed Points Are Not Empty:\\Hidden Internal Structure Behind Apparently Static Recurrent Surfaces}
\author{Author metadata to be supplied before submission}
\date{June 13, 2026}

\begin{document}

\maketitle

\begin{abstract}
Recurrent systems are often summarized by their exposed trajectory: fixed
point, cycle, transient, or chaotic regime. This paper argues that such
surface labels can be misleading. In Demian, a structured recurrent-substrate
research program, we repeatedly observe apparently static surface behavior
coexisting with differentiated internal state. Dual-GRU and native-substrate
experiments show fixed-point surfaces that split into distinct basin
interiors, including tight and accumulating fixed-point classes.
Capsule-continuity probes further show that full internal-state restore can
resume trajectories exactly while surface-only restore fails, implying that
continuation-relevant information is not contained in the exposed surface
alone. We compare these findings against plain RNN, GRU, LSTM, diagonal SSM,
Transformer-reservoir, and Mamba-reservoir controls, and we use negative
gate-state checks to bound mechanism claims. The result is not a claim of
benchmark superiority or a finished architecture, but a measurement protocol:
fixed-point behavior should be treated as a hypothesis about surface dynamics,
not as evidence that the internal computation is empty.
\end{abstract}

\section{Introduction}

Fixed points are often read as simple endpoints. In a recurrent system, a
state appears to settle; the exposed readout changes little; the trajectory is
classified as static. That label is useful, but it can also be destructive if
it is treated as a verdict about the computation inside the system.

The central claim of this paper is narrow: an apparently static surface can
hide structured internal state. A surface fixed point can coexist with
message-state accumulation, channel separation, route activity, and
continuation-relevant memory. Therefore a fixed-point label should be treated
as a surface classification, not as evidence that the internal computation is
empty.

This problem is not specific to one architecture. Prior work has used
dynamical-systems tools to reverse-engineer trained recurrent networks,
including fixed points and line-attractor structure
\citep{maheswaranathan2019lineattractor}. Other work shows that networks with
similar behavior can implement different internal solutions
\citep{huang2024degeneracy}, and that comparing dynamics requires tools beyond
static representation comparison \citep{guilhot2024dsa}. Adjacent work also
frames modern sequence models, including Transformers and state-space models,
as dynamical systems rather than only input-output maps
\citep{fernando2025transformerdynamics,gu2023mamba}.

Demian is a substrate-building research program whose laboratory artifacts
make this surface/internal split unusually explicit. The experiments studied
here include plain recurrent baselines, dual-GRU variants, native Demian
substrates, five-channel evolved substrates, Transformer self-reference
reservoirs, and Mamba self-reference reservoirs. The paper does not claim that
Demian beats these baselines. It uses them to make a measurement point: fixed
surface behavior is not enough.

The contribution is a compact evidence discipline for separating surface
collapse from internally structured fixed-point basins:

\begin{itemize}
  \item classify the exposed surface trajectory;
  \item separately classify fixed-point interiors;
  \item measure route and channel activity;
  \item perturb or ablate internal state;
  \item compare full internal-state resume against surface-only resume;
  \item demote mechanism names when strict held-out checks fail.
\end{itemize}

\section{Terminology}

We use \emph{surface} for the exposed readout vector of a recurrent system.
We use \emph{internal state} for the recurrent variables that are not fully
specified by that surface readout. A \emph{channel} is a named component of
internal state, and a \emph{route} is a learned or designed transformation
between state components.

A \emph{surface fixed point} is an exposed trajectory whose readout converges
or nearly converges under the classifier used by the experiment. An
\emph{interior fixed point} is a finer label describing the hidden state under
that surface class. The most important interior labels in this paper are:

\begin{description}
  \item[tight fixed point] a fixed surface with little measured internal
  accumulation under the selected metrics;
  \item[accumulating fixed point] a fixed surface with continuing internal
  accumulation or differentiated internal activity;
  \item[surface-only resume] a control that reconstructs only the exposed
  surface and zeros or omits hidden state;
  \item[full capsule resume] a control that restores the recurrent body and
  full internal state before continuing.
\end{description}

These terms are operational. They describe what the diagnostics measured in
the repository artifacts, not a claim of consciousness, agency, or global
architectural superiority.

\section{Methods}

\subsection{Substrates and baselines}

The evidence comes from five groups of runs:

\begin{enumerate}
  \item dual-GRU fixed-point interior maps;
  \item canonical Demian v9 and v9 five-channel capsule-continuity probes;
  \item matched plain recurrent and native-substrate perturbation baselines;
  \item Transformer and Mamba self-reference reservoir runs;
  \item gate-state truth-campaign checks used as negative controls.
\end{enumerate}

The artifact source is the Demian Lab repository \citep{demianlab2026}. The
paper's tracked evidence summary is included in
\path{papers/fixed_point_internal_structure/evidence_summary.csv}; the original
source artifacts remain in the repository's \texttt{data/} tree.

\subsection{Surface and interior measurements}

The first measurement layer labels the exposed trajectory as a fixed point,
cycle, transient, or related regime. The second layer asks whether fixed-point
surfaces are internally homogeneous. Dual-GRU summaries report interior
classes, message-state norms, bottleneck code counts, and bottleneck entropy.
These quantities are not presented as universal metrics; they are the local
measurement interface for this substrate family.

\subsection{Capsule-continuity probe}

Capsule continuity asks whether a paused trajectory can be resumed from
different reconstructions of state. Let $s_t$ be the internal recurrent state
and $y_t = R(s_t)$ be the exposed surface. The uninterrupted continuation is:
\[
  s_{t+k} = F^k(s_t).
\]
The full capsule control restores $s_t$ and compares its continuation with the
uninterrupted trajectory. The surface-only control constructs a state
$\hat{s}_t$ from $y_t$ alone, leaving the non-surface internal state empty or
zeroed, and compares:
\[
  F^k(\hat{s}_t) \quad \textrm{against} \quad F^k(s_t).
\]
If full capsule resume matches and surface-only resume diverges, the exposed
surface was not a sufficient state for continuation.

\subsection{Claim discipline}

The paper separates measurement claims from mechanism names. The gate-state
truth campaign is included because it demoted a stronger mechanism label under
stricter held-out checks. That negative result is part of the method: a
surface/internal difference is evidence for hidden structure, but not
automatically evidence for a named causal mechanism.

\section{Results}

\subsection{Fixed-point surfaces split into different interiors}

Table~\ref{tab:dual-gru} shows the central fixed-point interior result. Every
listed dual-GRU variant is classified as \texttt{FIXED\_POINT} in 8/8 runs at
the exposed surface. The internal readings are not identical. Some variants
split into both tight and accumulating fixed-point interiors, while
\texttt{dual\_gru\_v3b:current} is entirely accumulating and has much larger
mean message norm and bottleneck entropy.

\begin{table}[H]
\centering
\small
\begin{tabular}{llllrr}
\toprule
Substrate & Surface & Interior counts & Classes & Message norm & Entropy \\
\midrule
\texttt{dual\_gru\_v3b:current} & 8/8 fixed & accumulating: 8 & 1 & 23.1086 & 1.7779 \\
\texttt{dual\_gru\_v3b:tight} & 8/8 fixed & tight: 6; accumulating: 2 & 2 & 0.0624 & 0.3423 \\
\texttt{dual\_gru\_v3b:threshold} & 8/8 fixed & tight: 5; accumulating: 3 & 2 & 0.0753 & 0.5120 \\
\texttt{dual\_gru\_v3} & 8/8 fixed & tight: 6; accumulating: 2 & 2 & 0.0624 & 0.5069 \\
\bottomrule
\end{tabular}
\caption{Dual-GRU fixed-point surfaces with different basin interiors.
All rows share the same exposed surface attractor class, but differ in hidden
interior class, message-state expression, and bottleneck entropy.}
\label{tab:dual-gru}
\end{table}

This is the simplest evidence that a fixed-point label is not enough. The
surface class is constant across the table, while the internal class and
internal activity vary.

\subsection{Surface-only resume fails where full capsules resume}

Table~\ref{tab:capsule} shows capsule-continuity probes for canonical
\texttt{demian\_native\_v9} and the v9 five-channel substrate. Across the
reported sweep, full internal-state resume remains exact or near-exact, while
surface-only resume remains worse than full capsule resume.

\begin{table}[H]
\centering
\small
\begin{tabular}{lrrr}
\toprule
Substrate & Full capsule min cosine & Full capsule max gap & Surface-only min gap \\
\midrule
\texttt{demian\_native\_v9} & 0.99999988 & 0.0 & 0.22624873 \\
\texttt{v9\_five\_channel} & 0.99999988 & 0.0 & 0.27455845 \\
\bottomrule
\end{tabular}
\caption{Capsule-continuity summary. Full internal-state capsules resume
the deterministic continuation, while surface-only reconstruction does not.}
\label{tab:capsule}
\end{table}

This supports a direct operational claim: the visible surface is not a
sufficient state for continuation in these probes. The hidden state contains
continuation-relevant information.

\subsection{Matched baselines bound the interpretation}

Table~\ref{tab:baselines} summarizes matched pseudo-channel perturbation
baselines from the truth-campaign comparison. Plain RNN, GRU, and LSTM
baselines recover almost exactly under this protocol. Diagonal SSM and native
Demian substrates show larger perturbation consequences.

\begin{table}[H]
\centering
\small
\begin{tabular}{lrrr}
\toprule
Substrate & Runs & Final L2 gap mean & Recovery-window gap mean \\
\midrule
\texttt{rnn} & 3 & $1.19\times10^{-7}$ & 0.000111 \\
\texttt{gru} & 3 & $5.18\times10^{-8}$ & 0.000144 \\
\texttt{lstm} & 3 & $4.46\times10^{-8}$ & 0.000061 \\
\texttt{diag\_ssm} & 3 & 0.0650 & 0.002366 \\
\texttt{demian\_native\_v8} & 3 & 0.2206 & 0.003906 \\
\texttt{demian\_native\_v9} & 3 & 0.00322 & 0.000531 \\
\bottomrule
\end{tabular}
\caption{Matched perturbation/recovery baselines. The table is not a
superiority ranking. It shows why baseline controls are necessary before
promoting internal-structure claims.}
\label{tab:baselines}
\end{table}

The conservative reading is that some perturbation effects are ordinary
recurrent recovery behavior, while other effects require more careful
mechanism tests. This is why the paper centers measurement discipline rather
than global architecture claims.

\subsection{Architecture context: Transformer and Mamba reservoirs}

The Transformer and Mamba reservoir artifacts are useful context but not the
main proof. The Transformer self-reference reservoir supported a deterministic
period-2 attractor observation in the Demian claims ledger. Mamba
self-reference runs did not simply reproduce that period-2 behavior; in the
batch summary, cached Mamba runs have lag-2 autocorrelation 0.1495, while
no-cache runs show 570 period-2 switches in 1000 steps. This supports a
modest point: recurrence-like surface behavior depends on architecture, so
surface classes should not be treated as architecture-independent evidence.

\subsection{Negative gate-state checks keep the claim bounded}

The 2026-05-16 truth campaign demoted the stronger gate-state mechanism name
under strict held-out checks: the strict Track B profile pass rate was 0/2 in
the inspected multisurgery summary. This does not weaken the fixed-point
paper's main claim. It strengthens the paper's boundary. Hidden structure is
not the same as settled causal mechanism naming.

\section{Discussion}

The results support a practical distinction: fixed point is a surface class,
not a computational verdict. In the dual-GRU family, the same surface label
contains different interior classes. In the capsule-continuity probes, a
surface-only reconstruction fails to preserve continuation even when a full
internal capsule resumes exactly. In the matched baseline comparison, plain
recurrent models and native substrates behave differently under the same
pseudo-channel perturbation protocol.

The broader lesson is methodological. A recurrent system should not be judged
only by the final or exposed trajectory. Surface labels should be paired with
state-sensitive controls: interior class, route norms, channel separation,
ablation response, and resume fidelity.

This is also the design reason Demian v1 makes internal channels explicit. The
claim is not that v1 is complete. The claim is that the experimental record
made anonymous hidden state too weak an object of study. If continuation lives
inside the state, then the state must be named, saved, ablated, and tested.

\section{Limitations}

The evidence is heterogeneous because it comes from a living research
repository rather than a single polished benchmark suite. Some runs have small
seed counts. Some metrics are local to Demian's substrate families. The
Transformer and Mamba reservoir artifacts are architectural context rather
than central evidence. The capsule-continuity result uses full internal-state
restore, not a compressed state code. Finally, the paper does not prove that
every fixed point is meaningful; it proves that fixed-point surface labels can
be incomplete.

\section{Reproducibility}

The source repository is Demian Lab \citep{demianlab2026}. The primary
artifact files used in this draft are:

\begin{itemize}
  \item \path{data/substrate_lab/dual_gru_family_summary.json}
  \item \path{data/substrate_lab/v9_capsule_continuity_20260511/summary.json}
  \item gate-state truth campaign multisurgery baseline-comparison summary
  \item gate-state truth campaign multisurgery campaign summary
  \item \path{data/mamba_batch/mamba_batch_summary.json}
  \item \path{data/reservoir_batch/batch_summary.json}
\end{itemize}

The paper-local evidence table is stored beside this source as
\path{evidence_summary.csv} and \path{evidence_summary.json}. The paper can be
built with:

\begin{verbatim}
cd papers/fixed_point_internal_structure
make
\end{verbatim}

\section{Conclusion}

Fixed-point behavior is not empty by default. It is a statement about the
observed surface under a particular classifier. The Demian artifacts show that
fixed surfaces can hide accumulating interiors, differentiated state, and
continuation-relevant internal information. The right response is not to
promote every fixed point, but to stop treating the surface label as the whole
system.

\bibliographystyle{plainnat}
\bibliography{references}

\end{document}
